Pour une scintigraphie, on injecte du technétium 99m, qui émet un photon gamma de \qty{140}{keV} à chaque désintégration. Pendant \qty{\v{t}}{h}, l'activité dans un organe de \qty{\v{m}}{kg} est d'environ \qty{\v{Act}}{MBq} (on la suppose constante). L'organe absorbe \qty{\v{fr}}{\percent} de l'énergie des photons émis.

\begin{enumerate}[label=\alph*)]
    \item Combien de désintégrations se produisent dans l'organe pendant cette durée ?
    \item Quelle énergie l'organe absorbe-t-il ?
    \item Déterminez la dose absorbée et la dose équivalente ($w_R = 1$ pour les gammas), en millisieverts.
\end{enumerate}
